\documentclass[a4paper,11pt]{article}
\usepackage[T1]{fontenc}
\usepackage[utf8]{inputenc}
\usepackage[margin=2.5cm]{geometry}
\usepackage{amsmath,amssymb}
\usepackage{booktabs}
\usepackage{listings}
\usepackage{xcolor}
\usepackage{hyperref}

\lstset{basicstyle=\ttfamily\small,columns=fullflexible,keepspaces=true,mathescape=false,
  literate={⌊}{{$\lfloor$}}1 {⌋}{{$\rfloor$}}1 {₀}{{$_0$}}1 {₁}{{$_1$}}1 {≤}{{$\le$}}1}

\newcommand{\cuvee}{Cuvée}
\newcommand{\tr}[1]{\lfloor #1\rfloor}

\title{Cuvée: Verifying Loops by Trace Equation Systems\\and a Lemma Database\\[0.5ex]\large An overview of the
current approach}
\date{October 2026}

\begin{document}
\maketitle

\begin{abstract}
\cuvee{} verifies programs with loops without asking for loop invariants. Symbolic execution summarizes each loop
by a \emph{trace equation system}: the sequence of values that a group of variables takes, defined by its initial
values and the loop body as a fold over the iterations, together with the number of iterations and a constraint
that the loop test held in each of them. The goals that result are about such traces. A deduction engine rewrites
them with a database of lemmas about common loop shapes, counters, sums, searches, stores into arrays, and
recurrences, into first-order arithmetic that it decides itself, or hands what remains to an SMT solver. This
document describes the approach as it is implemented now, assuming familiarity with program verification but not
with trace equations, and collects the insights that shaped it.
\end{abstract}

\section{The problem and the idea}

Verifying a program with loops usually means finding inductive invariants, which is where automated verifiers
spend their effort: model checkers search for them, Horn clause solvers synthesize them, deductive verifiers ask the
user. \cuvee{} avoids invariants altogether. It describes what a loop \emph{computes}, not what it preserves.

The observation behind it is that most loops in practice are instances of a small number of recurring shapes: a
counter that is incremented, a sum or a count accumulated over the elements of an array, an array initialized
position by position, a flag set when a condition is met, a search that stops at the first match. For each shape,
the final values have a closed form or a simple characterization, e.g.\ the sum of the elements read, or the fact
that no element before the exit matched. If the effect of a loop is written down as an equation over the sequence
of its states, these characterizations become \emph{lemmas about sequences}, which can be proved once, collected in
a database, and applied by rewriting.

The approach therefore has three parts:
\begin{enumerate}
\item A translation of programs into \emph{trace equation systems} (TEQs), which is exact: no abstraction is involved,
  every loop becomes an equation that defines its traces, and the verification condition is a formula over these
  traces (Section~\ref{sec:teq}).
\item A \emph{lemma database} of facts about traces of common shapes, written as Dafny lemmas and proved in Lean
  (Section~\ref{sec:database}).
\item A \emph{deduction engine} that rewrites the verification conditions with the lemmas and decides the remaining
  arithmetic, by congruence closure and a simplex (Section~\ref{sec:engine}).
\end{enumerate}
What the engine cannot close is either lowered into first-order logic for an SMT solver or reported as unknown;
\cuvee{} does not need to know whether a goal is provable before trying.

\section{Traces and combinators}
\label{sec:traces}

A \emph{trace} is the non-empty sequence of values that a loop variable takes, one per iteration plus the initial
value. Four combinators describe traces and sequences derived from them; all take a list of \emph{dependencies},
sequences of equal length that are traversed point-wise:
\begin{itemize}
\item $\mathit{trace}(x_0,\ \textbf{for}\ Y \ \textbf{do}\ x, y :: u)$ starts with $x_0$ and appends, for each
  element $y$ of the dependency $Y$, the value of the update $u$ for the previous value $x$; it is a left fold that
  keeps all intermediate results. With several accumulators, $\mathit{trace}(x_0, z_0,\ \textbf{for}\ Y\
  \textbf{do}\ x, z, y :: u_x, u_z)$ folds a tuple, read by components $\mathit{current}(S).i$ and histories
  $\tr{S}.i$.
\item $\mathit{maps}(\textbf{for}\ Y\ \textbf{yield}\ y :: e)$ is the point-wise image, $\mathit{invs}(\textbf{for}\
  Y\ \textbf{check}\ y :: \varphi)$ the statement that $\varphi$ holds at every position.
\item $\mathit{current}(X)$ is the last value of a trace, $\tr{X}$ (``butlast'') all earlier ones, a list with one
  element per iteration, and $\mathit{earlier}(j, \tr{X})$ the value before iteration $j$.
\item $\mathit{counter}(i, c, n)$ is the trace $i, i + c, \dots, i + n c$ of $n$ steps; the iterations of every loop
  are counted by a counter $K = \mathit{counter}(0, 1, n)$, whose $n$ is the number of iterations.
\end{itemize}
An update may consist of guarded cases, $\textbf{case}\ \varphi_1 \Rightarrow e_1\ \textbf{case}\ \varphi_2
\Rightarrow e_2$, one per path through the loop body. The guards are disjoint by construction (they are the branch
conditions of the paths); they need not be complete, e.g.\ after a \texttt{break}, and a step where no guard holds
yields a value that is deliberately unspecified.

\section{From programs to trace equation systems}
\label{sec:teq}

\cuvee{} reads C (the SV-COMP dialect), Dafny, Boogie, SMT-LIB and SV-LIB into one core language. The transformation
\texttt{-teq} executes a program symbolically, inlining calls from the entry point (\texttt{main} in C), and treats
each loop on a single path:
\begin{itemize}
\item The variables that the loop modifies are partitioned by the \emph{strongly connected components} of their
  read relation: a variable reads the variables in its updates and in the guards of the paths. Each component
  becomes one trace: a variable that depends on no cycle gets a trace of its own, as in the original per-variable
  formulation; variables that depend on each other, e.g.\ the bounds $\mathit{lo}$ and $\mathit{hi}$ of a binary
  search, form one trace of several accumulators. The components are ordered so that each depends only on earlier
  ones, so no trace is defined in terms of a later one, and the system is \emph{acyclic}.
\item The number of iterations $n$ and the values of nondeterministic choices are fresh; the counter, the traces and
  the exit values $\mathit{current}(X)$ are defined by equations; the \emph{iteration constraint} states, by
  $\mathit{invs}$ over the counter and all histories, that the loop test held before each iteration; the exit
  states that it fails afterwards. The body is checked to execute safely from the final state.
\item Nested loops are closed: an inner loop becomes a value of the outer update, with its own counter and traces
  substituted by their definitions, rather than a family of sequences indexed by the outer iteration.
\item Exits by \texttt{break} or \texttt{return} are told apart by branch conditions, not by assumptions, so the
  cases of an enclosing loop stay disjoint.
\end{itemize}
The result is exact: the goal is valid iff the program satisfies its assertions (up to the modelled C semantics,
e.g.\ bounded integers with wrap-around for unsigned types). For the loop
\begin{lstlisting}
i := 0; s := 0; while (i < n) { s := s + a[i]; i := i + 1; }
\end{lstlisting}
\texttt{-teq} produces, with the counter $K = \mathit{counter}(0, 1, m)$,
\[ I = \mathit{counter}(0, 1, m),\quad S = \mathit{trace}(0,\ \textbf{for}\ \tr{I}\ \textbf{do}\ s, i :: s +
a[i]),\quad \mathit{invs}(\textbf{for}\ \tr{I}\ \textbf{check}\ i :: i < n),\quad \neg(\mathit{current}(I) < n), \]
and the lemma database turns $\mathit{current}(S)$ into $\mathit{sum}(\mathit{selects}(a, \tr{I}))$, the sum of
the elements read, and the iteration constraint with the exit into $\mathit{current}(I) = n$ for $0 \le n$.

\section{The lemma database}
\label{sec:database}

The database is a set of Dafny files: \texttt{teq.dfy}, always loaded, declares the core (naturals, lists, traces,
counters, constant traces, stores and reads of arrays) with the lemmas every loop needs, and optional modules add
concepts on demand: \emph{aggregation} (sums, counts, maxima, minima, parities), \emph{arithmetic} (powers, triangle
numbers, division), \emph{arrays} (stores at counters and keys, conditional stores, resets), \emph{counters} (phases,
flags, catching up, bounded values), \emph{search} (agreeing prefixes, absence of a value, bisection),
\emph{sorted} (sortedness, permutations, partition points), \emph{strides} (strided stores), and
\emph{unsigned} (counters modulo $2^{32}$). A lemma is used in one of two ways:
\begin{itemize}
\item A \emph{rewrite rule} (\texttt{\{:rw\}}) replaces its left-hand side, e.g.\ the trace of a loop that adds
  $c$ in each step is the counter, $\mathit{trace}(i,\ \textbf{for}\ \tr{K}\ \textbf{do}\ x, k :: x + c) =
  \mathit{counter}(i, c, n)$, and the final value of a counter is $i + n c$. Rules that \emph{name} a trace (fold it
  into a function of the database such as $\mathit{sums}$, $\mathit{stores}$ or $\mathit{bisection}$) come first,
  rules that \emph{compute} with named traces after them.
\item A \emph{forward lemma} (\texttt{\{:fw\}}) adds a fact wherever a term matching its trigger occurs among the
  assumptions, e.g.\ that a count is non-negative, or that a bisection keeps its lower bound below the cut of its
  guard. Facts that are needed where only a term occurs, as a bound for the arithmetic, belong here rather than in a
  rule about an atom that the goal may never contain.
\end{itemize}
Side conditions of rules are decided by the engine in the context of the goal, or looked up among its facts.

The lemmas are proved in Lean~4 against a semantics of the combinators (traces as non-empty lists, combinators as
folds); the statements are generated from the Dafny files, so the proofs are about exactly the lemmas that
\cuvee{} uses. A small number of premises that the engine does not check (reads within the range of the loop, equal
lengths) are listed explicitly. The proofs found real inconsistencies (Section~\ref{sec:insights}).

\section{The deduction engine}
\label{sec:engine}

The engine (\texttt{-prove}) simplifies a goal under the facts of its context, applying the rules where they match
and deciding atoms by its theories:
\begin{itemize}
\item \emph{Congruence closure} over the terms of the facts, which also provides the definitions of variables
  (members of their class) through which rules match: a rule on $\mathit{current}(\mathit{trace}(\dots))$ matches
  $\mathit{current}(S)$ if $S$ is defined as such a trace.
\item A \emph{simplex} over the rationals with integer tightening of bounds, which also knows the bounds of
  division by literals ($k (t\ \mathrm{div}\ k) \le t < k (t\ \mathrm{div}\ k) + k$).
\item \emph{Matching modulo} associativity and commutativity of sums, higher-order patterns in Miller's fragment
  ($F(k)$ matches any expression in $k$, by abstraction), the cases of updates (a case of the expression must be
  covered by a case of the pattern), and the dependencies of combinators (a selection, where rules written for one
  trace per variable also match the single dependency of a component).
\item \emph{Views} of traces with several accumulators: a read of one component is matched as the trace of that
  accumulator alone, the others bound over their histories, so lemmas about single traces apply to cycles; lemmas
  may also match the whole cycle, e.g.\ the bisection of a binary search.
\item \emph{Forward facts}, canonical forms of terms with binders (so that renamed or reordered conditions are
  recognized), and memos of the contexts of combinators.
\end{itemize}
Rewriting is bounded by a step budget per goal; a goal that exhausts it is reported as not terminating, which is
a non-answer like ``unknown''.

\section{Beyond rewriting}

Three further modes complete the tool. \texttt{-fallback} lowers what remains of a goal into first-order logic,
traces as uninterpreted functions of an index with their characterizing equations, naturals as non-negative
integers, and hands it to Z3; this proves many goals that need arithmetic or quantifier reasoning beyond the engine,
and refutes goals without traces. \texttt{-chc} translates the program into constrained Horn clauses instead, a
complementary technique for loops over integers. \texttt{-verify} answers \emph{correct}, \emph{incorrect} or
\emph{unknown} per verification task; a refutation counts only if the goal was not over-approximated on the way
(goals that the lowering generalized are marked abstract).

\section{Evaluation}

All runs: BenchExec, 15\,min CPU time, 7\,GB, 2 cores of an Intel Core i7-6700; the other tools are the official
results of SV-COMP 2026. On ReachSafety-Arrays (440 tasks) and ReachSafety-Loops (764):

\begin{center}\small
\begin{tabular}{l|rrr|rrr|r}
\toprule
& \multicolumn{3}{c|}{Arrays} & \multicolumn{3}{c|}{Loops} & \\
& true & false & wrong & true & false & wrong & score \\
\midrule
\cuvee{} (no solver) & 138 & 9 & 0 & 149 & 16 & 0 & 599 \\
\cuvee{} with Z3 & 148 & 48 & 0 & 230 & 20 & 0 & 824 \\
\cuvee{} with Horn clauses & 8 & 45 & 0 & 150 & 41 & 0 & 402 \\
portfolio of the last two & 151 & 67 & 0 & 262 & 52 & 0 & 945 \\
\midrule
Symbiotic (SV-COMP 2026) & & & & & & & 1099 \\
UAutomizer (SV-COMP 2026) & & & & & & & 941 \\
\bottomrule
\end{tabular}
\end{center}

The strength is arrays of unbounded size: the next tool proves 68 safe Arrays tasks, and 112 of the safe Arrays
tasks that \cuvee{} proves are proved by none of five SV-COMP~2026 participants; these need invariants quantified
over array indices. Loops over integers with non-linear invariants (\texttt{nla-digbench}) are the weakness, where
Horn clause solvers and model checkers lead; the solver helps there most (63 of the 83 proofs it adds on Loops).
TEQ and Horn clauses are complementary: their portfolio scores higher than either. On the MonoCera benchmark of
aggregates over arrays (125 tests of universal properties, maxima, minima and sums), \cuvee{} proves 79 and 88
with Z3, against 41 for MonoCera, a tool built for this class. No configuration gives a wrong answer.

\section{Insights}
\label{sec:insights}

\paragraph{Representation.}
\begin{itemize}
\item \emph{One trace per strongly connected component} is the right unit: per-variable traces of variables that
  depend on each other are defined through each other, which rules cannot match; a trace of several accumulators is
  acyclic, its components are recovered as views on demand, and lemmas can match the cycle as a whole.
\item Closing nested loops (an inner loop as a value of the outer update) removes a whole module of lemmas about
  fragments of inner loops that the per-variable form needed.
\item Exits must be told apart by branch conditions: as assumptions, two exits with different values gave
  overlapping cases, from which any goal followed (a wrong proof found on SV-COMP).
\end{itemize}

\paragraph{Soundness.}
\begin{itemize}
\item Unspecified values are a trap: a lemma that bounds a trace for any number of steps also bounds steps after
  the loop would have stopped, where the value is unspecified, and two instances can then derive
  $\mathit{false}$. The Lean proofs found this for the bisection; the fix is a premise on the rule that names the
  trace, that every step was taken with a guard, which the iteration constraint of the loop provides.
\item A refutation is only a counterexample if the goal was exact; every approximation must be tracked (abstract
  goals), and uninterpreted functions (bit operations) make satisfiability meaningless.
\item Checking procedures out of the context of their calls gave wrong refutations; the entry point fixes the
  context.
\end{itemize}

\paragraph{The database.}
\begin{itemize}
\item Order matters: rules are tried in the order of the files, so rules that name traces must precede rules that
  compute with them, and rules whose patterns subsume others (a sum $v + w$ also matches $-e$ as $-e + 0$) must come
  after the specific ones, or the matcher must exclude the degenerate match.
\item Facts about terms belong in forward lemmas, triggered by a \emph{named} term (a function of the database)
  rather than a raw trace pattern, which would be matched at every read of every component.
\item Dropping rules that look specific is risky: a migration of the database that removed rules never applied in a
  sample lost dozens of SV-COMP tasks (counters modulo $2^{32}$, sentinels); the usage statistics of full runs, not
  of samples, decide.
\item Many rules are still written for one trace per variable (iterations and loop variable as two dependencies);
  teaching the matcher to map them onto the single dependency of a component recovered them without rewriting them.
\end{itemize}

\paragraph{The engine.}
\begin{itemize}
\item Matching is where proofs fail silently: normal forms differ between rules and goals (negative literals,
  products without a constant, implied guards), and a rule that should apply does not. Each such gap, once found,
  closes whole families of tasks.
\item Proving \emph{safe} tasks is the strength of the approach and scores double; refuting unsafe tasks needs a
  model, which the engine does not build, and is left to the solver.
\item Performance matters only where it costs results: re-matching forward lemmas at every fact and re-simplifying
  siblings in each other's context are the main costs.
\end{itemize}

\section{Limitations}

Non-linear integer arithmetic, loops whose summary needs a recurrence without a closed form, arrays written at
data-dependent positions (pairwise cancelling sums), sorting algorithms with nested loops, and heap data
structures are beyond the current database and engine. Witnesses are not produced, so SV-COMP scores the verdicts
as unconfirmed. The engine has no integer reasoning about remainders beyond the bounds of division.

\end{document}
